\chapter{Полный гессиан вихря и барьер относительной жёсткости}

\section{Какой оператор требуется}

Предыдущий нерадиальный аудит проверил эффективную подсистему заряженной
амплитуды, нейтральной амплитуды и двух поперечных компонент связности.
Для окончательного решения нужны все пять компонент
$Q\in\operatorname{Sym}_0^2(V_1)$, все семь компонент
$T\in\operatorname{Sym}_0^3(V_1)$ и полная семейная связность.

Строгий журнал проекта, однако, содержит незакрытое условие. В главе о
масштабе и портале допустимое лоренцево продолжение было записано с
коэффициентами $Z_Q$ и $Z_F$, причём прямо установлено, что они не
следуют из конечного родителя. Последующие главы вывели кинетическую
норму поля $T$ из условного следа и коэффициент кривизны
\begin{equation}
 G=\frac{2}{27},
\end{equation}
но не вывели относительную пространственную жёсткость поля $Q$.

Это согласуется с литературной границей. Полная теория Ландау--де Жена
для $Q$-тензора в общем случае содержит несколько независимых упругих
инвариантов; см. \path{arXiv:2607.19920}, \path{arXiv:1910.01424} и
\path{arXiv:1906.09232}. Следовательно, выбор одного положительного числа
не является автоматической сменой единиц. Для спектра вихря также нужна
фоновая калибровка, отделяющая физические моды от чистых калибровочных;
см. \path{arXiv:1605.09175} и \path{arXiv:1602.09084}.

\section{Точный представительный состав}

Пусть $J$ генерирует вращения вокруг выбранной тетраэдрической оси.
Действие $-J^2$ на двух полях имеет спектры
\begin{align}
 \Spec(-J^2|_{V_2})&=\{0,1,1,4,4\},\\
 \Spec(-J^2|_{V_3})&=\{0,1,1,4,4,9,9\}.
\end{align}
Тем самым все внутренние угловые веса известны без выбора базиса.

Проекция канонического тетраэдра на ядро $J$ и его ортогональное
дополнение точно воспроизводит разложение профиля
\begin{equation}
 T_*=T_0+T_3,
 \qquad
 \|T_0\|^2=\frac{160}{81},
 \qquad
 \|T_3\|^2=\frac{128}{81},
 \qquad
 \langle T_0,T_3\rangle=0.
\end{equation}
Это независимо проверяет коэффициенты $A$ и $B$ предыдущей краевой
задачи.

\section{Максимальный однозначный радиальный оператор}

До вывода $Z_Q$ можно построить не один оператор, а семейство. В
совместно вращающемся осесимметричном секторе оно действует на
\begin{equation}
 5+7+1=13
\end{equation}
вещественных компонентах: полном $Q$, полном $T$ и угловой вариации
связности. Его квадратичная форма содержит
\begin{equation}
 Z_Q\|D_r\delta Q\|^2
 +\|D_r\delta T\|^2
 +\frac{G}{r^2}(\delta K')^2
 +\langle(\delta Q,\delta T),H_V(r)(\delta Q,\delta T)\rangle,
\end{equation}
а угловая часть определяется указанными выше весами $J^2$.

Локальный потенциальный гессиан вдоль профиля не является положительным:
его минимальная кривизна достигает
\begin{equation}
 \lambda_{\min}^{\mathrm{loc}}=-0.97583411\ldots.
\end{equation}
Поэтому устойчивость нельзя вывести из однородного вакуума; она зависит
от полного баланса производных и связности.

На сетке из $150$ радиальных узлов проверены значения
\begin{equation}
 Z_Q\in\{0.03,0.1,0.3,1,3,10,30\}.
\end{equation}
В каждом случае шесть нижних собственных значений положительны. Наименьшее
равно
\begin{equation}
 \lambda_{\min}=0.01466229243\ldots
\end{equation}
и в пределах скана практически не зависит от $Z_Q$. Следовательно, в
проверенном полном внутреннем радиальном секторе отрицательной моды не
обнаружено.

\section{Почему это ещё не полный гессиан}

Положительный скан не устраняет три разрыва:
\begin{enumerate}
 \item родитель не выбирает единственное значение $Z_Q$ и не исключает
 дополнительные упругие инварианты $Q$;
 \item не выведена единая временная кинетическая метрика, поэтому
 статические собственные значения нельзя считать физическими частотами;
 \item текущий радиальный оператор не содержит все поперечные компоненты
 неабелевой связности и не проверяет все неосесимметричные сектора.
\end{enumerate}

\begin{nogocriterion}[Барьер определённости полного гессиана]
Нельзя объявлять полную устойчивость вихря, пока один и тот же родительский
след не фиксирует относительную жёсткость $Q$, кинетику $T$, кривизну
семейной связности и временную норму. Скан произвольного положительного
$Z_Q$ является проверкой робастности, но не выводом физического оператора.
\end{nogocriterion}

Нулевая гипотеза следующего гейта: в текущей конечной архитектуре нет
ковариантного квадрата, который порождает пространственную норму $DQ$ с
нужной относительной нормировкой. Если это подтвердится, ветвь устойчивой
частицы требует нового родительского принципа и не может продолжаться к
хопфовой петле автоматически.

\begin{protocol}
\begin{enumerate}
 \item все внутренние компоненты $Q$ и $T$: \textbf{включены};
 \item представительные угловые веса: \textbf{выведены точно};
 \item отрицательная локальная кривизна ядра: \textbf{найдена};
 \item отрицательная радиальная мода при $0.03\le Z_Q\le30$:
 \textbf{не найдена};
 \item относительная жёсткость $Q$: \textbf{не выведена};
 \item полный неабелев и неосесимметричный оператор: \textbf{не замкнут};
 \item полная устойчивость вихря: \textbf{не доказана};
 \item следующий гейт: \textbf{родительская нормировка жёсткости $Q$}.
\end{enumerate}
\end{protocol}

Машинный сертификат сохранён в
\path{s2t_v6_bosonic_defect_q_tetrahedral_vortex_full_hessian_gate_results.json}.